\documentclass[../main.tex]{subfiles}
\ifSubfilesClassLoaded{\setcounter{chapter}{6}}{}
\begin{document}

\chapter{Supervised vs Unsupervised; \texttt{train\_test\_split}; Alur \texttt{fit}/\texttt{predict}}

\begin{subcpmk}
  \item \textbf{Sub-CPMK-P5:} Memecah data latih/uji; melatih estimator \sklearn{} dengan pola API konsisten (RPS Mg.\ 6).
\end{subcpmk}

\SesiRPS{6}{P5}{$\sim$170 menit}{CPMK-P3}

\section{Kemampuan akhir sesi}

Menjelaskan beda supervised/unsupervised; menggunakan \texttt{train\_test\_split} dengan \texttt{random\_state}; memanggil \texttt{fit}, \texttt{predict}, \texttt{score} pada data uji; mengenali gejala overfitting konsep.

\section{Teori}

\begin{konsep}
Estimator \sklearn{} memiliki antarmuka seragam: \texttt{fit(X, y)} untuk supervised dengan label; \texttt{predict} atau \texttt{predict\_proba}; \texttt{transform} untuk prapemrosesan. Data uji tidak boleh dipakai untuk \texttt{fit} estimator atau scaler.
\end{konsep}

\section{Contoh kode}

\begin{lstlisting}[caption=Split dan baseline DummyClassifier]
from sklearn.model_selection import train_test_split
from sklearn.dummy import DummyClassifier
from sklearn.datasets import load_iris

X, y = load_iris(return_X_y=True)
X_train, X_test, y_train, y_test = train_test_split(
    X, y, test_size=0.2, random_state=42, stratify=y
)
base = DummyClassifier(strategy="most_frequent")
base.fit(X_train, y_train)
print("akurasi baseline", base.score(X_test, y_test))
\end{lstlisting}

\begin{lstlisting}[caption=LogisticRegression sederhana]
from sklearn.linear_model import LogisticRegression
clf = LogisticRegression(max_iter=200)
clf.fit(X_train, y_train)
print("akurasi model", clf.score(X_test, y_test))
\end{lstlisting}

\begin{lstlisting}[caption=Prediksi dan bentuk array]
import numpy as np
pred = clf.predict(X_test)
print(pred[:5], y_test[:5])
\end{lstlisting}

\begin{lstlisting}[caption=Catatan Markdown di notebook]
# Setelah split, semua fit/scaling hanya pada X_train
# X_test hanya untuk evaluasi akhir
\end{lstlisting}

\section{Referensi}

\begin{itemize}[leftmargin=*]
  \item sklearn basic tutorial. \url{https://scikit-learn.org/stable/tutorial/basic/tutorial.html}
  \item Kaggle Intro ML. \url{https://www.kaggle.com/learn/intro-to-machine-learning}
\end{itemize}

\begin{asesmen}
\textbf{Kuis singkat:} definisi overfitting satu kalimat + baris split terlihat di notebook.
\end{asesmen}

\begin{rangkuman}
Minggu 6 menstandarkan alur latih/uji dan API \sklearn{}.

\textbf{Kata kunci:} \texttt{train\_test\_split}, \texttt{fit}, \texttt{predict}
\end{rangkuman}

\end{document}
